\section{Linear Algebra}
\label{sec:linalg}
\lipsum[1-2]

\subsection{Vector spaces}
\lipsum[3-4]
\begin{tcolorbox}[TUMBox=TUMGreen, title=Definition, breakable]
  Sei $K$ ein Körper. Eine Menge $V$ zusammen mit zwei Verknüpfungen
  \begin{equation}
    +: V \times V \rightarrow V, (a,b) \rightarrow a+b \quad\text{und}\quad
    \cdot: K \times V \rightarrow V, (\lambda,b) \rightarrow \lambda\cdot b
  \end{equation}
  heißen $\boldmath{K}$\textbf{-Vectorraum} wenn gilt:\par

  \begin{description}[itemsep=0pt]
    \item[(V1)] Die Verknüpfung $+$ ist assoziativ, das heißt:
      \begin{equation}
        (x+y)+c = x+(y+z) \quad\forall x,y,z \in V.
      \end{equation}
    \item[(V2)] Die Verknüpfung $+$ ist kommutativ, das heißt:
      \begin{equation}
        x+y=y+x \quad\forall x,y \in V.
      \end{equation}
    \item[(V3)] Es gibt ein neutrales Element bezüglich $+$, genannt $0_V$
      (Nullvector), das heißt:
      \begin{equation}
        0_V+v=v+0_V=v \quad\forall v \in V.
      \end{equation}
    \item[(V4)] Jedes Element $v \in V$ besitz ein inverses Element bezüglich
      $+$, geschrieben $-v$, das heißt:
      \begin{equation}
        v+(-v)=-v+v=0 \quad\forall v \in V.
      \end{equation}
    \item[(V5)] Die Verknüpfung $\cdot$ ist assoziativ, das heißt:
      \begin{equation}
        \lambda\cdot(\mu\cdot v) = (\lambda\mu)\cdot v \quad\forall\lambda,\mu
        \in K \forall v \in V.
      \end{equation}
    \item[(V6)] Es gibt ein neutrales Element bezüglich $\cdot$, genannt das
      Einselement $1$ des Körpers $K$, das heipt:
      \begin{equation}
        1\cdot v = v\cdot 1 \quad\forall v \in V.
      \end{equation}
    \item[(V7)] Es gilt das Distributivgesetz:
      \begin{equation}
        \lambda\cdot(v+w) = \lambda\cdot v + \lambda\cdot w \quad\forall
        \lambda\in K, \forall v \in V.
      \end{equation}
  \end{description}
  Die Elemente $v \in V$nennen wir nun \textbf{Vektoren}. Die Verknüpfung
  $+$ heißt \textbf{Vektoraddition} und $\cdot$ heißt
  \textbf{Skalarmultiplikation}.
\end{tcolorbox}
\lipsum[5-6]
\begin{TUMBoxFill}[TUMOrange]{Alert}
  This is an improtant fact to remember!
\end{TUMBoxFill}
\lipsum[7]

\subsection{Matrices and linear transformations}
\lipsum[8-12]

\subsection{Eigenvectors and Eigenvalues}
\lipsum[13-20]

\section{Analysis}
\label{sec:analysis}
\lipsum[21-22]

\subsection{Differentials}
\lipsum[23-25]
\begin{tcolorbox}[TUMBox, title=Of utmost importance]
  \lipsum[1]
  \begin{equation}
    \frac{d}{dx}x^n = nx^{n-1}
  \end{equation}
  \lipsum[2]
\end{tcolorbox}
\lipsum[26-29]

\subsection{Integrals}
\lipsum[30-35]

\section*{Summary}
\label{sec:summary}
\lipsum[36]
\begin{figure}[tbp]
  \centering
  \includegraphics[width=.4\textwidth]{example-image}
  \caption{Another very important figure.}
\end{figure}
\lipsum[37-40]

\section*{Exercises}
\label{sec:exercises}
\lipsum[41-43]
